\documentclass[12pt]{article}
\usepackage[utf8]{inputenc}
\usepackage[T1]{fontenc}
\usepackage{hyperref}
\title{Evenki Language Resources for K-12 Curriculum Development}
\author{Compiled by Claude Sonnet 5.5 (claude-sonnet-5-5) for Andr\'as Kornai}
\date{}
\begin{document}
\maketitle

\section{Introduction}
Evenki (ISO 639-3 \texttt{evn}) is a Tungusic language spoken in Russia, China and Mongolia. The classification given for this survey concerns China, where no standard orthography exists; texts appear in a Latin pinyin-like alphabet, in Mongolian script, or in Chinese characters. This report rests on one round of searches and is short.

\section{Primary Educational Resources}
A 2012 blog post reports that Evenki primary and middle school textbooks have been published, but not clearly on the Chinese side; no catalogue of them was found: \url{https://bruce-humes.com/2012/05/31/tibetan-daur-evenki-and-oroqen-databases-appraised-by-experts/}

\section{Linguistic and Reference Materials}
\begin{itemize}
\item Evenki language (Wikipedia): \url{https://en.wikipedia.org/wiki/Evenki_language}
\item Evenki orthography (Wikipedia; mentions the 1998 Evenki-Chinese dictionary \emph{Ewengki nihang bilehu biteg}): \url{https://en.wikipedia.org/wiki/Evenki_orthography}
\item Nedjalkov, \emph{Evenki} (Routledge descriptive grammar): \url{https://www.routledge.com/Evenki/Nedjalkov/p/book/9780415600057}
\item Endangered Languages Project entry: \url{https://www.endangeredlanguages.com/lang/evn}
\end{itemize}

\section{Religious and Community Resources}
None found.

\section{Digital and Audio Resources}
Discussion of Tibetan, Daur, Evenki and Oroqen databases (2012): \url{https://lists.osu.edu/pipermail/mclc/2012-May/000695.html}

\section{Curriculum Development Context}
No current Chinese Evenki school curriculum was identified. Chinese library catalogues and regional publishers would need to be consulted directly.
\end{document}